% ------------------------------------------------------------
% Page 101
% ------------------------------------------------------------

J’encaisse la monnaie et prends la porte sans un mot. A-t-il compris que je m’adressais au
potard ?

Le Ricou m’avait bien prévenu. Je n’avais rien obtenu, avec, cependant, un petit espoir, vite
déçu, pour la fenêtre brisée. Après la politesse obligée du début, n’obtenant que des refus,
j’étais devenu progressivement agressif, irrespectueux, animé par la colère des timides, et
enfin méprisant. Je ne regrettais rien, j’avais libéré tout ce que j’avais sur le cœur.

\textbf{Le Pasteur}

C’est dans cette humeur-là que j’arrivais au presbytre. Là, surprise ! Je connais bien le Pasteur,
et surtout sa jeune femme native du Pompidou ; Ils sont à Meyrueis depuis deux ans. Je suis
très aimablement reçu, et nous évoquons beaucoup de souvenirs communs du Pompidou,
avant d’aborder le problème de la fenêtre. J’évoque le refus du Maire qui se déclare non-
concerné. Le Pasteur me précise que la paroisse est fort désargentée, et qu’il ne peut engager
ladite paroisse pour satisfaire à ma demande, pourtant financièrement fort modeste.

A demi-mot je crois comprendre qu’entre la paroisse, propriétaire du Temple et la mairie,
utilisatrice de l’école, les obligations de part et d’autre sont mal définies ; le Pasteur n’a pas
connaissance d’un quelconque contrat qui aurait dû être établi. S’agit-il d’un accord tacite, en
application de la loi Falloux, en vertu de laquelle les instituteurs dits communaux étaient plus
ou moins en obédience religieuse dans les communes où coexistaient plusieurs communautés,
comme c’était le cas au Pompidou au moins jusqu’à 1863 ?

Je compris la prudence du Pasteur au nom de la paroisse propriétaire du Temple, qui, en
acceptant de remplacer la fenêtre créerait un précédent dont la Mairie pourrait user pour se
dégager de toute obligation de réparation ou d’entretien autre que celle concernant la seule
salle de classe. Tout cela paraissait bien compliqué dans cette affaire de droit entre propriétaire
et locataire.

Le Pasteur me confirma que, deux ou trois fois, à la belle saison, il célébrait le culte aux
Oubrets, devant une bonne assistance. Je l’invitais avec son épouse, à déjeuner ce qu’ils
acceptèrent volontiers.

\textbf{L’Institutueur : M. Galtier}

Je pris congé, désireux de contacter l’instituteur nouvellement nommé au chef-lieu de canton.
Monsieur Galtier me reçut comme on reçoit un jeune collègue. Je lui fis part de mes démêlés
avec le Maire Il n’en fut pas surpris :

« Mes rapports avec lui ne sont guère aimables, mais j’ai ici l’avantage d’être soutenu par un
noyau laïque très actif, qui n’hésite pas à intervenir. Le débat est permanent, bien sûr teinté de
politique. L’animateur est un jeune professeur très militant ». \textit{(Il s’agit de M. Verdier qui sera
d’ailleurs élu député de Paris à la Libération).}

M. Galtier me remonta un peu le moral :
« Neuf mois, ça passe vite même aux Oubrets ».
